\documentclass[../DoAn.tex]{subfiles}
\begin{document}

Chương này trình bày các công nghệ được sử dụng trong WoodCert Auction, vai trò của chúng đối với yêu cầu ở Chương~\ref{chapter:Related_works} và lý do lựa chọn. Phiên bản được đối chiếu từ tệp khai báo phụ thuộc và cấu hình hiện hành; nội dung chương không được xem là kết quả đo hiệu năng của hệ thống.

\section{Tổng quan lựa chọn công nghệ}

WoodCert Auction sử dụng kiến trúc nguyên khối phân mô-đun (modular monolith) cho phần mềm phía máy chủ (backend) và ứng dụng trang đơn (Single Page Application -- SPA) cho phần mềm giao diện (frontend). Phía máy chủ chạy trong một tiến trình Spring Boot nhưng mã nguồn được phân chia theo miền nghiệp vụ. Một số luồng liên mô-đun giao tiếp qua cổng hoặc giao diện lập trình, song các phụ thuộc trực tiếp vẫn tồn tại; vì vậy, hệ thống không phải kiến trúc vi dịch vụ (microservices).

Kiến trúc này phù hợp với phạm vi đồ án vì chỉ cần một đơn vị triển khai chính, vẫn hỗ trợ giao dịch cơ sở dữ liệu giữa các chức năng liên quan và có chi phí vận hành phù hợp với một máy chủ ảo.

\begin{figure}[htbp]
    \centering
    \includegraphics[width=0.92\linewidth,keepaspectratio]{chapter-3/system_architecture.pdf}
    \caption{Kiến trúc triển khai của hệ thống WoodCert Auction}
    \label{fig:chapter3-system-architecture}
    \caption*{\textit{Nguồn: Tổng hợp từ cấu hình và mã nguồn hệ thống.}}
\end{figure}

Hình~\ref{fig:chapter3-system-architecture} thể hiện Nginx tiếp nhận HTTPS và chuyển tiếp yêu cầu tới các cổng loopback của hai ứng dụng. MySQL và Redis chỉ giao tiếp trong mạng Docker Compose, không công bố cổng ra Internet. Trình duyệt tải tệp trực tiếp lên Cloudinary bằng thông tin có chữ ký tạm thời. Với VNPay Sandbox, trình duyệt được chuyển hướng tới trang thanh toán và quay lại điểm cuối Return của hệ thống sau khi hoàn tất.

\section{Nền tảng phía máy chủ và bảo mật}

\subsection{Java 17 và Spring Boot 3.5.15}

Phần mềm phía máy chủ sử dụng Java 17 và Spring Boot 3.5.15. Spring Boot cung cấp cơ chế cấu hình, tiêm phụ thuộc và tích hợp các thành phần trong hệ sinh thái Spring \cite{walls2022spring,springdocs}. Đồ án sử dụng Spring Web cho REST API, Spring Data JPA cho dữ liệu quan hệ, Spring Validation cho dữ liệu đầu vào, Spring Mail cho email định danh, Spring WebSocket cho sự kiện thời gian thực và Spring Actuator cho điểm kiểm tra tình trạng dịch vụ.

Các tác vụ nền theo lịch sử dụng \texttt{@Scheduled} để kích hoạt và đóng phiên đấu giá, xử lý quá hạn thanh toán, quá hạn xác nhận gửi hoặc bàn giao hàng và tự động hoàn tất giao nhận theo điều kiện cấu hình. Giao dịch cơ sở dữ liệu được xác định bằng \texttt{@Transactional}; các luồng có cạnh tranh cập nhật còn kết hợp khóa bản ghi, khóa lạc quan (optimistic locking), khóa thao tác nghiệp vụ lũy đẳng và ràng buộc duy nhất.

\subsection{Spring Security và JWT HS512}

JSON Web Token (JWT) là định dạng biểu diễn tập thông tin xác nhận dưới dạng JSON, được chuẩn hóa trong RFC 7519 \cite{rfc7519}. WoodCert Auction sử dụng mã truy cập ký bằng HMAC-SHA-512, có thời gian sống 15 phút và chỉ được giữ trong bộ nhớ của phần mềm giao diện. Sau khi xác minh chữ ký và thời hạn, hệ thống tiếp tục kiểm tra trạng thái tài khoản và truy vấn quyền hiệu lực từ MySQL.

Mã làm mới có thời gian sống 7 ngày, được đặt trong cookie \texttt{HttpOnly}, \texttt{Secure}, \texttt{SameSite=Lax} và giới hạn đường dẫn \texttt{/api/v1/auth}. Phía máy chủ chỉ lưu giá trị băm, luân chuyển mã sau mỗi lần sử dụng và thu hồi khi đăng xuất hoặc đặt lại mật khẩu. Vòng đời mã làm mới được quản lý bền vững trong MySQL.

Spring Security tắt bộ lọc CSRF mặc định cho API. Hai thao tác sử dụng cookie là làm mới phiên và đăng xuất được bảo vệ bằng cơ chế cookie gửi kép (double-submit cookie) ở tầng ứng dụng. Redis ghi số lần đăng nhập thất bại theo email và chặn tạm thời trong 15 phút sau 5 lần thất bại.

\section{Dữ liệu, trạng thái thời gian thực và tích hợp ngoài}

\subsection{Flyway và MySQL 8.0}

Flyway quản lý thay đổi lược đồ bằng các tệp di trú có phiên bản và lưu lịch sử thực thi trong cơ sở dữ liệu \cite{flywaydocs}. Cấu trúc hiện hành được tạo từ năm tệp di trú, từ \path{V1__baseline_schema.sql} đến \path{V5__add_shipment_deadline.sql}. Cách quản lý này giúp bảng, khóa ngoại, chỉ mục, dữ liệu tham chiếu và hạn xác nhận gửi hàng được áp dụng theo cùng thứ tự giữa các môi trường.

MySQL 8.0 lưu dữ liệu quan hệ bền vững gồm tài khoản, sản phẩm, thẩm định, đấu giá, tài chính, đơn hàng, giao nhận, tranh chấp và nhật ký quản trị \cite{mysqldocs}. Tệp đa phương tiện nằm tại Cloudinary; MySQL chỉ lưu siêu dữ liệu (metadata) và quan hệ nghiệp vụ. InnoDB cung cấp giao dịch, khóa và ràng buộc khóa ngoại.

\subsection{Redis 7.4, Lua và WebSocket}

Redis 7.4 lưu trạng thái thời gian chạy của phiên \texttt{ACTIVE} bằng cấu trúc Hash và Set \cite{redisdocs}. Dữ liệu gồm giá, người dẫn đầu, thời điểm kết thúc và danh sách người đủ điều kiện. Đoạn mã Lua kiểm tra thời gian, tư cách tham gia, người dẫn đầu, bước giá và điều kiện gia hạn cuối phiên trong một lần thực thi nguyên tử. Cấu hình vận hành thực tế bật nhật ký ghi nối tiếp (Append Only File -- AOF) để hỗ trợ phục hồi sau khi Redis khởi động lại.

Lượt đặt giá được gửi qua REST API. Sau khi Redis chấp nhận, phía máy chủ phát sự kiện \texttt{NEW\_BID} qua WebSocket/STOMP rồi cố gắng lưu lịch sử và bản chụp trạng thái sang MySQL nhưng không bảo đảm bước ghi thứ cấp luôn thành công. Hệ thống cũng cố gắng lưu các lượt hết thời gian hoặc không đạt bước giá để phục vụ truy vết. Bộ môi giới thông điệp đơn giản (simple broker) được nhúng trong Spring, không phải dịch vụ độc lập.

WebSocket cung cấp kênh truyền hai chiều trên kết nối lâu dài theo RFC 6455 \cite{rfc6455}. STOMP tổ chức đăng ký nhận sự kiện theo chủ đề (topic), còn SockJS cung cấp phương án dự phòng khi không thiết lập được WebSocket trực tiếp \cite{springwebsocketdocs}. Kênh hiện hành chỉ phát sự kiện tới trình duyệt theo dõi; lệnh đặt giá không được gửi qua STOMP.

Bảng~\ref{tab:chapter3-storage-responsibility} phân biệt dữ liệu và trách nhiệm của MySQL, Redis cùng tác vụ nền trong vòng đời phiên đấu giá.

\begin{table}[htbp]
    \centering
    \caption{Phân chia trách nhiệm giữa MySQL, Redis và tác vụ nền}
    \label{tab:chapter3-storage-responsibility}
    \small
    \renewcommand{\arraystretch}{1.2}
    \begin{tabular}{|>{\raggedright\arraybackslash}p{0.16\linewidth}|>{\raggedright\arraybackslash}p{0.30\linewidth}|>{\raggedright\arraybackslash}p{0.42\linewidth}|}
        \hline
        \textbf{Thành phần} & \textbf{Dữ liệu chính} & \textbf{Trách nhiệm} \\ \hline
        MySQL & Cấu hình phiên, người tham gia, lịch sử, bản chụp trạng thái và trạng thái kết thúc & Lưu dữ liệu bền vững, ràng buộc quan hệ, phục vụ truy vấn, quyết toán và xử lý sau phiên \\ \hline
        Redis & Giá, người dẫn đầu, thời điểm kết thúc và danh sách người đủ điều kiện của phiên \texttt{ACTIVE} & Quyết định và cập nhật lượt đặt giá nguyên tử bằng Lua trong thời gian phiên hoạt động \\ \hline
        Tác vụ nền & Mã phiên đến hạn và trạng thái xử lý còn dang dở & Kích hoạt, chốt phiên, quyết toán tiền cọc, tạo đơn hàng và sửa các trường hợp chưa hoàn tất \\ \hline
    \end{tabular}
\end{table}

Khi kích hoạt phiên, tác vụ nền gọi tiến trình xử lý vòng đời để khóa bản ghi, đọc người tham gia có cọc \texttt{FROZEN}, khởi tạo Redis và chuyển phiên sang \texttt{ACTIVE}. Khi chốt phiên, tiến trình ưu tiên trạng thái cuối trong Redis; nếu khóa không còn dữ liệu, hệ thống khôi phục từ bản chụp và lịch sử hợp lệ trong MySQL. Trạng thái kết thúc được ghi bền vững trước khi quyết toán tiền cọc và tạo đơn hàng.

\subsection{Cloudinary, VNPay Sandbox và SMTP}

Cloudinary lưu trữ, biến đổi và phân phối tệp đa phương tiện qua mạng phân phối nội dung \cite{cloudinarydocs}. Phía máy chủ tạo thông tin tải lên có chữ ký; trình duyệt tải tệp trực tiếp rồi yêu cầu hệ thống xác nhận loại và kích thước. Các loại đã tích hợp gồm ảnh đại diện, ảnh sản phẩm, ảnh chứng minh thẩm định và bằng chứng tranh chấp. Hệ thống đã khai báo kỹ thuật cho video đóng gói nhưng chưa tích hợp loại tệp này vào luồng giao nhận.

VNPay Sandbox được dùng cho luồng nạp tiền thử nghiệm. Phía máy chủ tạo URL thanh toán có chữ ký; điểm cuối Return xác nhận kết quả sau khi kiểm tra mã kiểm tra (checksum), mã đơn vị, số tiền và trạng thái giao dịch. SMTP gửi email xác thực tài khoản và đặt lại mật khẩu bằng cấu hình từ biến môi trường.

\section{Công nghệ phía giao diện}

\subsection{React 19.2.6, TypeScript 5.9.3 và Vite 7.3.3}

React xây dựng giao diện theo thành phần và cơ chế kết xuất khai báo \cite{reactdocs}. TypeScript bổ sung kiểm tra kiểu tĩnh cho JavaScript \cite{bierman2014typescript}; Vite cung cấp máy chủ phát triển, thay thế mô-đun nhanh và quy trình đóng gói ứng dụng \cite{vitedocs}. Ba công nghệ này được dùng để xây dựng SPA cho khu vực công khai và các khu vực chức năng dành cho người tham gia/người mua, người bán, thẩm định viên và quản trị viên.

Các tuyến giao diện được tải trì hoãn, còn Vite chia gói theo nhóm phụ thuộc để hạn chế dữ liệu phải tải ở lần truy cập đầu. TypeScript phát hiện sai lệch kiểu khi xây dựng, nhưng dữ liệu từ API vẫn cần được kiểm tra và xử lý lỗi ở thời gian chạy.

\subsection{Quản lý dữ liệu, trạng thái và giao diện}

TanStack Query 5.100.9 quản lý dữ liệu máy chủ, bộ nhớ đệm và vòng đời yêu cầu trong ứng dụng React \cite{reactquerydocs}. Truy vấn đọc dữ liệu mặc định thử lại một lần khi lỗi, trong khi thao tác thay đổi dữ liệu không tự động thử lại nhằm hạn chế lặp các thao tác tài chính hoặc chuyển trạng thái. Bộ nhớ đệm được xóa khi đăng xuất hoặc khi danh tính người dùng thay đổi.

Zustand 5.0.13 giữ mã truy cập và trạng thái xác thực cục bộ trong bộ nhớ \cite{zustanddocs}. Tailwind CSS 4.2.4 cung cấp các lớp tiện ích để xây dựng bố cục \cite{tailwinddocs}; Radix UI cung cấp các thành phần nguyên thủy chú trọng hành vi tương tác và khả năng tiếp cận \cite{radixdocs}. Đồ án chỉ sử dụng một số thành phần Radix cần thiết, không xem đây là bộ giao diện hoàn chỉnh.

\section{Hạ tầng triển khai, CI/CD và kiểm thử}

\subsection{Nginx, Docker và Docker Compose}

Nginx trên máy chủ đóng vai trò proxy ngược và điểm kết thúc TLS; Nginx trong container giao diện phục vụ tệp tĩnh \cite{nginxdocs}. Các đường dẫn \texttt{/api/} và \texttt{/ws-auction} được chuyển tới phần mềm phía máy chủ, còn các đường dẫn khác tới phần mềm giao diện. Chứng chỉ HTTPS được thiết lập bằng Certbot và Let's Encrypt \cite{certbotdocs}.

Docker đóng gói hai ứng dụng thành các ảnh Docker (Docker image) riêng; Docker Compose mô tả MySQL, Redis, mạng nội bộ, vùng lưu trữ, kiểm tra tình trạng và thứ tự khởi động \cite{dockerdocs,composedocs}. MySQL và Redis không công bố cổng ra máy chủ; hai ứng dụng chỉ liên kết cổng vào loopback để lưu lượng Internet đi qua Nginx.

\subsection{GitHub Actions và kiểm thử tự động}

Quy trình tích hợp liên tục chạy kiểm thử Maven; kiểm tra quy tắc mã nguồn, kiểu dữ liệu, Vitest và đóng gói phần mềm giao diện; chạy Playwright; kiểm tra script triển khai, Docker Compose và việc xây dựng ảnh Docker \cite{actionsdocs}. Khi CI trên nhánh chính thành công, quy trình phát hành tạo ảnh gắn với mã commit và đẩy lên GitHub Container Registry.

Tác vụ triển khai chỉ chạy khi biến \texttt{CD\_ENABLED} bằng \texttt{true} và người có quyền phê duyệt GitHub Environment \texttt{production}. Sau khi được duyệt, quy trình chuyển script lên VPS qua SCP, thực thi qua SSH và kiểm tra điểm cuối sẵn sàng (readiness). Bí mật ứng dụng nằm trong tệp cấu hình được bảo vệ trên VPS, không được đóng gói vào ảnh Docker hoặc lưu trong kho mã nguồn.

Phía máy chủ sử dụng JUnit Jupiter 5.12.2 và Testcontainers để kiểm thử logic, cấu hình, MySQL, Redis và tệp di trú \cite{junitdocs,testcontainersdocs}. Phần mềm giao diện sử dụng Vitest, Testing Library và Playwright. Các công cụ này hỗ trợ kiểm thử logic, thành phần và luồng trình duyệt đầu cuối \cite{vitestdocs,playwrightdocs}. Kết quả phải dựa trên báo cáo của lần chạy tương ứng; mục tiêu hiệu năng và khả dụng chỉ được công nhận khi có đo kiểm hoặc dữ liệu giám sát có thể tái lập.

\section{Đối chiếu lựa chọn công nghệ}

Bảng~\ref{tab:chapter3-technology-comparison} tổng hợp các phương án thay thế chính. Việc lựa chọn dựa trên yêu cầu nghiệp vụ, trạng thái mã nguồn và chi phí triển khai của đồ án.

\begin{table}[H]
    \centering
    \caption{Đối chiếu công nghệ được lựa chọn và phương án thay thế}
    \label{tab:chapter3-technology-comparison}
    \footnotesize
    \renewcommand{\arraystretch}{1.15}
    \begin{tabular}{|>{\raggedright\arraybackslash}p{0.14\linewidth}|>{\raggedright\arraybackslash}p{0.18\linewidth}|>{\raggedright\arraybackslash}p{0.20\linewidth}|>{\raggedright\arraybackslash}p{0.31\linewidth}|}
        \hline
        \textbf{Nhóm} & \textbf{Lựa chọn} & \textbf{Phương án thay thế} & \textbf{Lý do phù hợp với đồ án} \\ \hline
        Phía máy chủ & Spring Boot & NestJS, Django & Phù hợp mã nguồn Java, giao dịch quan hệ và hệ sinh thái bảo mật, dữ liệu, kiểm thử \\ \hline
        Phía giao diện & React & Vue, Angular & Phù hợp cấu trúc SPA, thành phần và bộ kiểm thử hiện có \\ \hline
        Dữ liệu bền vững & MySQL & PostgreSQL & Lược đồ, tệp di trú, lớp truy cập dữ liệu và kiểm thử đã được xây dựng trên MySQL \\ \hline
        Trạng thái phiên & Redis và Lua & MySQL với khóa, Hazelcast & Hỗ trợ cấu trúc dữ liệu cần thiết và thực thi kiểm tra, cập nhật nguyên tử \\ \hline
        Truyền sự kiện & WebSocket và STOMP & SSE, thăm dò HTTP & Phù hợp kết nối lâu dài và mô hình chủ đề dùng chung cho các màn hình theo dõi \\ \hline
        Tệp đa phương tiện & Cloudinary & Amazon S3, lưu trên VPS & Hỗ trợ tải trực tiếp có chữ ký, đọc metadata và phân phối qua CDN \\ \hline
        Triển khai & Docker Compose và GitHub Actions & Triển khai thủ công, Kubernetes & Phù hợp một VPS, hỗ trợ kiểm tra tự động, ảnh bất biến và quy trình phê duyệt \\ \hline
        Kiểm thử & JUnit, Vitest, Playwright, Testcontainers & Các framework tương đương & Khớp hệ sinh thái Java, React và nhu cầu chạy dịch vụ phụ thuộc trong container \\ \hline
    \end{tabular}
\end{table}

\section{Tổng kết chương}

Chương~\ref{chapter:Methodology} đã trình bày các nền tảng tạo nên WoodCert Auction và lý do lựa chọn trong phạm vi đồ án. Spring Boot và MySQL quản lý logic cùng dữ liệu bền vững; Redis và Lua xử lý trạng thái phiên \texttt{ACTIVE}; WebSocket/STOMP phân phối sự kiện; React và các thư viện liên quan tổ chức SPA; Cloudinary, VNPay Sandbox và SMTP hỗ trợ tích hợp ngoài; Docker, Nginx và GitHub Actions hỗ trợ triển khai có kiểm soát. Các công nghệ này là cơ sở cho phần phân tích thiết kế, triển khai và đánh giá ở Chương~\ref{chapter:Experiment}.

\FloatBarrier
\end{document}
